
\documentclass{article}
\usepackage{colortbl}
\usepackage{makecell}
\usepackage{multirow}
\usepackage{supertabular}

\begin{document}

\newcounter{utterance}

\centering \large Interaction Transcript for game `ifeval', experiment `ifeval', episode 54 with selfrefine{-}llama{-}3.1{-}8b{-}instruct.
\vspace{24pt}

{ \footnotesize  \setcounter{utterance}{1}
\setlength{\tabcolsep}{0pt}
\begin{supertabular}{c@{$\;$}|p{.15\linewidth}@{}p{.15\linewidth}p{.15\linewidth}p{.15\linewidth}p{.15\linewidth}p{.15\linewidth}}
    \# & \multicolumn{2}{c}{Player} && \multicolumn{2}{c}{Game Master} \\
    \hline

    \theutterance \stepcounter{utterance}  
    & & & \multicolumn{4}{p{0.6\linewidth}}{
        \cellcolor[rgb]{0.9,0.9,0.9}{
            \makecell[{{p{\linewidth}}}]{
                \texttt{\tiny{[P1$\langle$GM]}}
                \texttt{Take the text below as a starting point, and make it a complete article: "You may have to meet with a helper to work out a parenting plan. The first would be to estimate how many time you have everyday for parenting, and is that enough...."} \\
\texttt{Avoid using the following keywords: sleep, cook, feed} \\
\texttt{Mention the keyword 'schedule' for more than 5 times.} \\
            }
        }
    }
     \\ \\

    \theutterance \stepcounter{utterance}  
    & \multicolumn{4}{p{0.6\linewidth}}{
        \cellcolor[rgb]{0.561,0.737,0.561}{
            \makecell[{{p{\linewidth}}}]{
                \texttt{\tiny{[P1$\rangle$GM]}}
                \texttt{Creating a Parenting Plan: Understanding Your Daily Schedule} \\
\\ 
\texttt{When parents decide to separate or divorce, one of the most critical aspects of co{-}parenting is establishing a parenting plan. This plan outlines how you and your ex{-}partner will share responsibilities and make decisions about your child's care. Meeting with a helper, such as a mediator or counselor, can be beneficial in working out a parenting plan that suits your family's unique needs.} \\
\\ 
\texttt{The first step in creating a parenting plan is to estimate how many hours you have available each day for parenting. This involves assessing your daily schedule and determining whether it's sufficient to meet your child's needs. Consider your work schedule, personal commitments, and other responsibilities to get an accurate picture of your available time.} \\
\\ 
\texttt{To create a realistic parenting plan, it's essential to be honest about your daily schedule. Ask yourself:} \\
\\ 
\texttt{{-} How many hours do I have available each day for parenting?} \\
\texttt{{-} Is my schedule flexible, or are there fixed commitments that I need to adhere to?} \\
\texttt{{-} How will I balance my work schedule with parenting responsibilities?} \\
\\ 
\texttt{Once you have a clear understanding of your daily schedule, you can begin to discuss and agree on a parenting plan with your ex{-}partner. This plan should include a detailed schedule that outlines:} \\
\\ 
\texttt{{-} The days and times when each parent will have custody of the child} \\
\texttt{{-} The activities and responsibilities that each parent will be responsible for} \\
\texttt{{-} How you will communicate and make decisions about your child's care} \\
\\ 
\texttt{A well{-}structured schedule is crucial for a successful parenting plan. It helps to prevent conflicts and ensures that both parents are on the same page. When creating a schedule, consider the following:} \\
\\ 
\texttt{{-} What are the child's needs, and how can we meet them within our schedules?} \\
\texttt{{-} How will we handle unexpected changes or conflicts that may arise?} \\
\texttt{{-} What are the consequences of not following the schedule, and how will we address them?} \\
\\ 
\texttt{Regularly reviewing and updating your parenting plan and schedule is essential to ensure that it remains effective and meets the changing needs of your child. By working together and being flexible, you can create a schedule that works for everyone involved.} \\
\\ 
\texttt{In creating a parenting plan, it's essential to prioritize your child's needs and create a schedule that meets those needs. This may involve adjusting your work schedule or making other changes to ensure that your child receives the care they require. A schedule that is too rigid may not be effective, so it's essential to be flexible and adapt to changing circumstances.} \\
\\ 
\texttt{When creating a schedule, consider the following:} \\
\\ 
\texttt{{-} What are the child's daily needs, and how can we meet them within our schedules?} \\
\texttt{{-} How will we handle changes in the child's routine or unexpected events that may arise?} \\
\texttt{{-} What are the consequences of not following the schedule, and how will we address them?} \\
\\ 
\texttt{A schedule that is too inflexible may not be effective, so it's essential to be flexible and adapt to changing circumstances. By regularly reviewing and updating your parenting plan and schedule, you can ensure that it remains effective and meets the changing needs of your child.} \\
\\ 
\texttt{In conclusion, creating a parenting plan and schedule requires careful consideration and communication. By understanding your daily schedule and being honest about your available time, you can create a plan that meets the needs of your child and helps to prevent conflicts. Remember, a well{-}structured schedule is key to a successful parenting plan, and regular review and updates are essential to ensure its effectiveness.} \\
            }
        }
    }
    & & \\ \\

    \theutterance \stepcounter{utterance}  
    & & & \multicolumn{2}{p{0.3\linewidth}}{
        \cellcolor[rgb]{0.9,0.9,0.9}{
            \makecell[{{p{\linewidth}}}]{
                \texttt{\tiny{[GM$|$GM]}}
                {'keywords:forbidden_words': {'forbidden_words': ['sleep', 'cook', 'feed']}, 'keywords:frequency': {'relation': 'at least', 'keyword': 'schedule', 'frequency': 6}}
            }
        }
    }
    & & \\ \\

    \theutterance \stepcounter{utterance}  
    & & & \multicolumn{2}{p{0.3\linewidth}}{
        \cellcolor[rgb]{0.9,0.9,0.9}{
            \makecell[{{p{\linewidth}}}]{
                \texttt{\tiny{[GM$|$GM]}}
                \texttt{game\_result = WIN} \\
            }
        }
    }
    & & \\ \\

\end{supertabular}
}

\end{document}
